\documentclass{article}
\usepackage[utf8]{inputenc}
\usepackage{amsmath,amssymb}
\usepackage[most]{tcolorbox}
\usepackage{geometry}
\geometry{margin=2.5cm}

\newcommand{\step}[1]{\par\noindent\hangindent=3.5em\hangafter=1%
  \makebox[3.5em][l]{\textbf{#1.}}}
\newcommand{\steptag}[1]{\hfill{\small\textsf{[#1]}}}

\newtcolorbox{proofbox}[1]{
  colback=gray!5, colframe=gray!60,
  fonttitle=\bfseries, title={#1},
  breakable, enhanced,
  left=4pt, right=4pt, top=4pt, bottom=4pt,
  before skip=10pt, after skip=10pt
}

\begin{document}

\begin{proofbox}{If B is a finite-dimensional C*-algebra, A is a finite-di...}
\step{1} If B is a finite-dimensional C*-algebra, A is a finite-dimensional extended epsilon-C*-algebra, v:B->A is linear, and 0 <= delta <= delta', then if v is an extended delta-inclusion it is an extended delta'-inclusion, and if v is an extended delta-isomorphism it is an extended delta'-isomorphism and in particular an extended delta'-inclusion. \steptag{validated} \\[6pt]
\step{1.1} Assume v is an extended delta-inclusion and fix arbitrary n>=1. By def-extended-delta-inclusion, v\_n=1\_\{M\_n\} tensor v is a delta-homomorphism and, for every X, (1-delta)||X|| <= ||v\_n(X)|| <= (1+delta)||X||. By GT-kitaev-def-delta-homomorphism, v\_n is bounded linear, satisfies ||v\_n(I)-I|| <= delta and ||v\_n(XY)-v\_n(X)v\_n(Y)|| <= delta||X||||Y||, and in the *-algebra setting preserves the involution exactly. Since 0<=delta<=delta' and norms and products of norms are nonnegative, these inequalities imply ||v\_n(I)-I|| <= delta' and ||v\_n(XY)-v\_n(X)v\_n(Y)|| <= delta'||X||||Y||; bounded linearity and exact involution preservation are unchanged. Thus v\_n is a delta'-homomorphism by GT-kitaev-def-delta-homomorphism. Moreover (1-delta')||X|| <= (1-delta)||X|| and (1+delta)||X|| <= (1+delta')||X||, since ||X||>=0, so v\_n satisfies the two-sided (1+-delta') norm bounds. As n was arbitrary, def-extended-delta-inclusion makes v an extended delta'-inclusion. \steptag{validated} \\[4pt]
\step{1.2} Assume v is an extended delta-isomorphism. By def-extended-delta-inclusion this means that v is an extended delta-inclusion and v is bijective. The first child makes the same map v an extended delta'-inclusion, while bijectivity is independent of the defect parameter and is unchanged. Hence the definition makes v an extended delta'-isomorphism; and every extended delta'-isomorphism is, by that same definition, in particular an extended delta'-inclusion. \steptag{validated} \\[4pt]
\end{proofbox}

\end{document}
